\section{Algoritmo di Ant Colony Optimization}
L'algoritmo di \textbf{A}nt \textbf{C}olony \textbf{O}ptimization (\textbf{ACO}) è un algoritmo probabilistico 
per identificare un cammino hamiltoniano all'interno di un grafo pesato. Per cammino hamiltoniano si intende un cammino che passa una e una sola volta per tutti i nodi di un grafo.
Questo algoritmo è basato sul comportamento delle formiche, le quali rilasciano dei feromoni per comunicare tra di loro.

Inizialmente descriviamo l'algoritmo da un punto di vista superficiale, entrando incrementalmente nel dettaglio.

Specifichiamo che l'algoritmo utilizzato è di tipo asincrono: i nodi ACO del grafo non hanno necessità di coordinarsi tra di loro per funzionare.

Per iniziare i nodi ACO hanno a disposizione un sottoinsieme di nodi del grafo e gestiscono l'inizializzazione e lo scambio di formiche presenti nei nodi ACO.
Una formica rappresenta un'agente che si muove all'interno del grafo per cercare il cammino migliore. 
Durante la scelta del cammino i nodi ACO devono calcolare la probabilità di scelta di un determinato arco. La probabilità dipende da 3 elementi: la quantità di feromoni depositata su quell'arco, dal peso associato all'arco e da una somma di tutti i feromoni e pesi associati a tutti gli archi disponibili. 
Una volta trovata una soluzione, cioè la formica riesce a tornare al punto di partenza, il nodo ACO aggiorna i feromoni associati agli archi che la formica ha visitato.

Quindi, più dettagliatamente l'algoritmo ACO segue delle regole che determinano:
\begin{itemize}
    \item la quantità di feromoni depositata negli archi attraversati dalle formiche;
    \item l'arco scelto dalla formica durante il suo cammino;
    \item la velocità con cui i feromoni evaporano.
\end{itemize}
Una formica posizionata nel nodo \texttt{i} decide di muoversi nel nodo \texttt{j} con probabilità $p_{i,j}$:
\[ p_{i,j} = \frac{{(\tau_{i,j})}^{\alpha}{(\eta_{i,j})}^\beta}{\sum_{j \in \delta(i,j)}{(\tau_{i,j})}^{\alpha}{(\eta_{i,j})}^\beta} \] 
dove:
\begin{itemize}
    \item $\tau_{i,j}$ è la quantità di feronome depositata nell'arco $(i,j)$;
    \item $\eta_{i,j}$ è l'inverso del peso associato all'arco $(i,j)$;
    \item $\alpha$ è un parametro per controllare l'influenza di $\tau_{i,j}$;
    \item $\beta$ è un parametro per controllare l'influenza di $\eta_{i,j}$;
    \item $\delta(i,j)$ è una lista di nodi raggiungibili dal nodo $i$ non ancora visitati.
\end{itemize}
Le soluzioni migliori devono essere ``marchiate'' da dei feromoni per indicare a future formiche il percorso ottimale trovato.
Di conseguenza, quando una formica $a$ completa un cammino $L_a$ di costo $C_a$, 
migliore dell'ultimo cammino conosciuto, le formiche incrementeranno la quantità di 
feromoni presenti in ogni arco del cammino con un valore pari a $\Delta\tau_{i,j}^a$ che è proporzionale alla qualità della soluzione:
\[ \Delta\tau_{i,j}^a = 
    \begin{cases}
    {\frac{1}{C_a}}       & \quad \text{if edge } (i,j) \in L_a\\
    0                     & \quad \text{otherwise}
    \end{cases}
\]
Ricordiamo che $L_a$ è il costo del cammino della formica $a$. Quando una formica viaggia lungo un cammino, viene tenuto conto della lunghezza del cammino. 

L'aggiornamento del valore dei feromoni tiene in considerazione anche l'evoporazione di essi, risultando:
\[
    \tau_{i,j} = (1 - \rho) \tau_{i,j} + \rho \Delta \tau_{i,j}^\alpha
    \]
dove: $\rho$ è il coefficiente di evaporazione dei feromoni e $0 \leq \rho < 1$.

Sorge però un problema, le formiche usando la formula di probabilità\dots
Aggiungere feromoni locali e manca la spiegazione per step dell'algoritmo.
Dopo aver definito tutto segue la lista puntata delle azioni da intraprendere.



